% !TEX root = Hauptdokument.tex

\usepackage[utf8]{inputenc}
\usepackage[ngerman]{babel}
\usepackage[left=2.5cm,right=5.5cm,top=4.5cm,bottom=2.4cm]{geometry}
\usepackage[T1]{fontenc}

\usepackage{amsmath}
\usepackage{amsfonts}
\usepackage{amssymb}

\usepackage{graphicx}
\usepackage{subfig} %Mehrere Bilder gemeinsame Unterschrift
\usepackage{pgfplots}
\usetikzlibrary{calc,patterns,decorations.pathmorphing,decorations.markings}
\usetikzlibrary{arrows.meta}
\pgfkeys{/pgf/number format/.cd ,use comma ,set thousands separator={.}} % Zahlen in pgfplots werden über die pgf Engine gesetzt. Siehe Doku zu pgfplotstable 
\pgfplotsset{ every axis label/.append style={font=\scriptsize}, % 
every tick label/.append style={font=\scriptsize}} % globale Einstellungen. 

% nur eine Säule anzeigen
\pgfplotsset{/pgfplots/ybar legend/.style={
		/pgfplots/legend image code/.code={\draw[##1,/tikz/.cd,bar
width=3pt,yshift=-0.2em,bar shift=0pt] plot coordinates {(0cm,0.8em)
%hier geloescht (2*\pgfplotbarwidth,0.6em)
};},
	}}

\usepackage{tikz}

\usepackage{fancyhdr}
\pagestyle{fancy}
\usepackage{multirow} % Tabelle mit multi Zeile
\usepackage{colortbl} % Tabellen einfärgen

\usepackage{color} %inkscape
\usepackage{transparent} %inkscape
\usepackage{floatflt} %Textfluss um Bilder 
\usepackage{float} %Textfluss um Bilder
\usepackage{wrapfig} %Textfluss um Bilder
\usepackage{epstopdf} %eps-Dateien verarbeiten
\usepackage{pdflscape} %Seitenorientierung
\usepackage{array} %Tabellenumgebung erweitern
\usepackage{tabularx} % umfassendere Tabellen
\usepackage{xcolor} %Farben definieren
\usepackage{setspace} %Zeilenabstände
\usepackage{acronym} %Abkürzungsverzeichnis, nur verwendete in Verzeichnis aufnehmen mit [printonlyused]
\usepackage[per=slash,
            decimalsymbol=comma,
            loctolang={DE:ngerman,UK:english},
            ]{siunitx} %SI-Einheiten richtig darstellen
\usepackage{nicefrac} %Schräger Bruchstrich
\sisetup{exponent-product = \cdot, output-product = \cdot} % \cdot anstelle x in wissenschaftl. Notation Siunitx Paket



% Formelzeichenverzeichnis
\newcommand{\acrounit}[1]{
  \acroextra{\makebox[18mm][l]{\si{#1}}}
}
\usepackage{array}
\usepackage{rotating}
\usepackage{ragged2e}
\newcolumntype{C}[1]{>{\centering}p{#1}} % neuer Spaltentyp Tabelle
\newcolumntype{L}[1]{>{\RaggedRight\arraybackslash}p{#1}} % Neue
                                % Tabellenspalte: linksbündig mit
                                % Breitenangabe

\pgfplotsset{compat=newest}
\usepgfplotslibrary{units}

%\usepgfplotslibrary{external} 
%\tikzexternalize[prefix=Plots/]
\usepackage{pdfpages} %zb includepdf u. Seitenorientierung

%%%%%%%%Schriftfarben%%%%%%%%%%%%%
\definecolor{Weinrot}{RGB}{139,0,0}
\definecolor{Weinrot2}{RGB}{178,34,34}
\definecolor{Hellgruen}{RGB}{144,238,144}
\definecolor{Limette}{RGB}{50,205,50}
\definecolor{Blau1}{RGB}{95,158,160}
\definecolor{Blau2}{RGB}{70,130,180}
\definecolor{Blau3}{RGB}{30,144,255} 

%HS-Esslingen Schriftfarben
\definecolor{HSBlau1}{RGB}{0,70,102} %Überschriften
\definecolor{HSBlau2}{RGB}{85,146,172} %Überschriften
\definecolor{HSGrau1}{RGB}{236,242,245} 
\definecolor{HSGrau2}{RGB}{245,248,249} %Hintergrund Abb.
\definecolor{HSGrau3}{RGB}{112,113,114} %Aufzählungszeich.
\definecolor{HSRot1}{RGB}{212,0,50} %Überschriften
\definecolor{HSRot2}{RGB}{230,143,126} 

%%%%%%%%%%%%%%%%%%%%%%%%%%%%%%%%%%

%%%%%%%%%%LAYOUT-EINSTELLUNGEN%%%%%%%%%%%

% Zitationsstil
\usepackage[style=authoryear,backend=biber,sorting=nyt]{biblatex}
\addbibresource{Literatur.bib}

%\addtokomafont{disposition}{\fontfamily{arial}} %Schriftart

\setcounter{secnumdepth}{4} %Tiefe Kapitelnummerierung

\setlength{\parindent}{0pt} %Einzug neuer Absatz

% Abstand zur Kopfzeile einstellen
\renewcommand{\chapterheadstartvskip}{\vspace{0cm}}

%%%KOPF-FUSSZEILEN
\fancypagestyle{plain} %Kopf u Fußzeilen auch auf Kapitelseiten

\fancyhead[L]{}
\fancyhead[C]{}
\fancyhead[R]{\includegraphics[trim= 0 0 0 0, clip,width=4cm]{HE-Logo.png}}
\fancyfoot[L]{}
\fancyfoot[C]{}
\fancyfoot[R]{\thepage}

\renewcommand{\headrulewidth}{0.5pt}
\renewcommand{\footrulewidth}{0.5pt}

